% Volume III: Attention and Publicity
% Part VI. Privacy as Boundary and Capacity
% Chapter 32: Privacy as Protection Against Model Capture
% Spine: proof-spines/ch032-spine.md

\chapter{Privacy as Protection Against Model Capture}
\label{ch:ch032}

\textbf{Model capture} names the condition in which an individual's actionable institutional representation becomes difficult to inspect, correct, or escape. Formally, the danger appears when
\[
M_t(x) \to A_t(x)
\]
\Definition\ when there is no adequate route from the represented person back to \(M_t\). Publicly available traces can then be pooled into a durable model that shapes access to employment, insurance, credit, mobility, reputation, or legal scrutiny before the subject can contest what pattern is being used.

The central harm is not observation alone, but the externalization and operationalization of a person-pattern as someone else's asset. Once cross-context traces are stabilized into a predictive representation, self-presentation becomes reactive and suspicion becomes effectively unilateral. Privacy therefore functions as protection against the capture, refinement, and deployment of these models. The same barriers to capture become concealment when institutions claim proprietary or security privacy for high-stakes models while still acting on them against the people they score.
